\documentclass[11pt,a4paper]{article}
\usepackage[margin=0.85in,top=1in,bottom=1in]{geometry}
\usepackage{amsmath,amssymb,amsfonts}
\usepackage{graphicx}
\usepackage{float}
\usepackage{enumitem}
\usepackage{microtype}
\usepackage{caption}
\usepackage[hidelinks]{hyperref}
\usepackage[most,breakable]{tcolorbox}
\tcbuselibrary{breakable,skins}
\usepackage[utf8]{inputenc}
\usepackage[T1]{fontenc}
\usepackage{lmodern}
\captionsetup{font=small,labelfont=bf,skip=4pt}
\setlist{leftmargin=1.5em,topsep=2pt}

%% Breakable card --- prevents content cutoff across pages
\newtcolorbox{solcard}[3]{
  breakable, enhanced,
  colback=white, colframe=gray!50, arc=2pt,
  title={\textbf{#1}\hfill\textsf{\small #3\,pts}},
  fonttitle=\normalsize\sffamily\bfseries,
  before upper={\small\textit{#2}\par\smallskip\hrule height 0.4pt\smallskip},
  top=4pt, bottom=6pt, left=7pt, right=7pt,
  parbox=false,
}

\title{\textbf{Synchrotron physics --- Solution}}
\author{\normalsize Y26 (2026) --- English translation}
\date{}
\begin{document}\maketitle\vspace{-0.5em}

\begin{solcard}{A1}{Determine the magnitude of the magnetic field $B(t)$ of the electromagnets at the moment when the energy of the electrons is $E(t)$. Express your answer in terms of $E(t)$, $T$ and physical constants.}{0.50}

Let us write down Newton's second law: $$e v B = p\,\frac{2\pi}{T}.$$ Then $$B = \frac{2\pi}{eT}\cdot\frac{p}{v} = \frac{2\pi}{eT}\cdot\gamma m = \frac{2\pi}{eT}\cdot\frac{E}{c^2}.$$

\par\smallskip\noindent\textbf{Answer:} $$B(t) = \frac{2\pi E(t)}{eT c^2}.$$
\end{solcard}

\begin{solcard}{A2}{Determine the relation between $T$, $U$ and $dB/dt$ for which equilibrium electrons exist.}{0.40}

Let us express the energy from the result of the previous point and differentiate, taking into account that $T = \text{const}$: $$\frac{dE}{dt} = \frac{eTc^2}{2\pi}\cdot\frac{dB}{dt} = \frac{eU}{T}\cos\varphi,$$ $$\cos\varphi = \frac{T^2c^2}{2\pi U}\cdot\frac{dB}{dt}.$$ The cosine does not exceed $1$, which gives the condition on the admissible parameters of the synchrotron.

\par\smallskip\noindent\textbf{Answer:} $$\frac{T^2c^2}{2\pi U}\cdot\frac{dB}{dt} \leqslant 1.$$
\end{solcard}

\begin{solcard}{A3}{Find the dependence of the frequency $\omega_U(t)$ on time. Express your answer in terms of $B(t)$, $n$, $R$ and physical constants.}{0.50}

Newton's second law is analogous to the previous points: $$e v B = p\,\frac{2\pi}{T} = \frac{p\,\omega_U}{n}.$$ Express $p$ and $E$: $$p = eB\,\frac{vT}{2\pi} = eBR,$$ $$E = \sqrt{(m_p c^2)^2 + (pc)^2} = \sqrt{(m_p c^2)^2 + (ceBR)^2}.$$ The relation between $E$ and $B$ can also be obtained analogously to the previous point: $$eB = \frac{p\,\omega_U}{vn} = \frac{E\,\omega_U}{c^2 n},$$ $$\omega_U(t) = \frac{nec^2 B}{E}.$$ It only remains to substitute the value of the energy.

\par\smallskip\noindent\textbf{Answer:} $$\omega_U(t) = \frac{nec^2 B(t)}{\sqrt{(m_p c^2)^2 + \left(ceB(t)R\right)^2}}.$$
\end{solcard}

\begin{solcard}{B1}{Obtain the value of $\alpha$ for the case of weak focusing.}{0.50}

Analogously to the previous point, the expression for the momentum on a circular orbit of radius $R$ is $$p(R) = eB(R)\,R = \text{const}\cdot R^{1-\xi}.$$ Then the relation between relative changes is $$\frac{dp}{p} = (1-\xi)\frac{dR}{R} = (1-\xi)\frac{dL}{L}.$$

\par\smallskip\noindent\textbf{Answer:} $$\alpha = \frac{1}{1-\xi}.$$
\end{solcard}

\begin{solcard}{B2}{Express $K$ in terms of $\gamma$ and $\alpha$.}{1.50}

Let us use the expression for the cyclic frequency: $$\omega = \frac{2\pi v}{L},$$ $$\frac{d\omega}{\omega} = \frac{dv}{v} - \frac{dL}{L} = \frac{dv}{v} - \alpha\frac{dp}{p},$$ $$K = \frac{E}{v}\frac{dv}{dE} - \alpha\frac{E}{p}\frac{dp}{dE}.$$ Let us find the derivatives and substitute: $$\frac{v}{E}\frac{dE}{dv} = \frac{v}{\sqrt{1-v^2/c^2}}\cdot\frac{d}{dv}\!\left(\frac{1}{\sqrt{1-v^2/c^2}}\right)\bigg/\frac{1}{\sqrt{1-v^2/c^2}} = \frac{v^2}{c^2-v^2} = \beta^2\gamma^2 = \left(1-\frac{1}{\gamma^2}\right)\gamma^2 = \gamma^2-1,$$ $$\frac{p}{E}\frac{dE}{dp} = \frac{p}{\sqrt{(mc)^2+p^2}}\frac{d\sqrt{(mc)^2+p^2}}{dp} = \frac{p^2}{(mc)^2+p^2} = 1-\frac{(mc^2)^2}{E^2} = 1-\frac{1}{\gamma^2} = \frac{\gamma^2-1}{\gamma^2}.$$

\par\smallskip\noindent\textbf{Answer:} $$K = \frac{1-\alpha\gamma^2}{\gamma^2-1}.$$
\end{solcard}

\begin{solcard}{B3}{Indicate for which type of focusing a critical energy exists. Find $\gamma$ for particles with the critical energy.}{0.50}

For the critical situation we have $$\gamma = \frac{1}{\sqrt{\alpha}}.$$ Since $\gamma > 1$, for the existence of a critical energy it is necessary that $\alpha < 1$. For weak focusing this condition is not met.

\par\smallskip\noindent\textbf{Answer:} A critical energy exists only for strong focusing, $$\gamma = \frac{1}{\sqrt{\alpha}}.$$
\end{solcard}

\begin{solcard}{C1}{Taking into account that the initial energies of the particles differ only slightly, express $\dot\varphi$ in terms of $K$, $\mathcal E$, $n$, $\omega_0$ and $E_0$.}{0.90}

Let us consider the change of the particle's phase relative to the equilibrium one over one revolution: $$\Delta\varphi = \omega_U\!\left(\frac{2\pi}{\omega} - \frac{2\pi}{\omega_0}\right).$$ Dividing by the time of this change, i.e.\ by the particle's period of revolution $2\pi/\omega$, we obtain $$\dot\varphi = \omega_U\!\left(1-\frac{\omega}{\omega_0}\right) = n(\omega_0 - \omega).$$ Using the definition of $K$, we obtain $$\omega - \omega_0 = K\,\frac{\omega_0}{E_0}\,\mathcal E.$$

\par\smallskip\noindent\textbf{Answer:} $$\dot\varphi = -\frac{nK\omega_0}{E_0}\,\mathcal E.$$
\end{solcard}

\begin{solcard}{C2}{Express $\dot{\mathcal E}$ in terms of $\varphi_0$, $\varphi$, $U$, $\omega_0$ and physical constants. Use the fact that the difference between $\omega$ and $\omega_0$ is small enough that it can be neglected here.}{0.30}

Let us write down the derivatives of $E$ and $E_0$: $$\dot E = \frac{\omega}{2\pi}\,eU\cos\varphi \approx \frac{\omega_0}{2\pi}\,eU\cos\varphi,$$ $$\dot E_0 = \frac{\omega_0}{2\pi}\,eU\cos\varphi_0.$$

\par\smallskip\noindent\textbf{Answer:} $$\dot{\mathcal E} = \frac{\omega_0\,eU}{2\pi}\left(\cos\varphi - \cos\varphi_0\right).$$
\end{solcard}

\begin{solcard}{C3}{Obtain an expression for $\ddot\varphi$ in terms of $K$, $E_0$, $\omega_0$, $U$, $\varphi_0$, $\varphi$, $n$ and physical constants. Determine for which $\varphi_0$ the equilibrium orbit is stable (that is, the initially infinitesimal deviations of $\varphi$ do not grow with time).}{0.60}

Let us differentiate $\dot\varphi$ taking into account the indicated approximations: $$\ddot\varphi = -\frac{nK\omega_0}{E_0}\,\dot{\mathcal E}.$$

For stability, expanding to first order: $$\ddot\varphi = \frac{nK\omega_0^2\,eU}{2\pi E_0}\sin\varphi_0\,(\varphi - \varphi_0).$$ Stability requires $K\sin\varphi_0 < 0$.

\par\smallskip\noindent\textbf{Answer:} $$\ddot\varphi = -\frac{nK\omega_0^2\,eU}{2\pi E_0}\left(\cos\varphi - \cos\varphi_0\right);$$ stability when $K\sin\varphi_0 < 0$.
\end{solcard}

\begin{solcard}{C4}{Express $\dot\varphi^2$ in terms of $K$, $E_0$, $\omega_0$, $U$, $\varphi_0$, $\varphi$, $n$, $\varphi_i$, $\mathcal E_i$ and physical constants.}{0.50}

Let us use the result of the previous paragraph: $$\ddot\varphi = \frac{d\dot\varphi}{dt} = \frac{d\dot\varphi}{d\varphi}\frac{d\varphi}{dt} = \frac{1}{2}\frac{d\dot\varphi^2}{d\varphi},$$ $$d\dot\varphi^2 = -\frac{nK\omega_0^2\,eU}{\pi E_0}\left(\cos\varphi - \cos\varphi_0\right)d\varphi,$$ $$\dot\varphi^2 + \frac{nK\omega_0^2\,eU}{\pi E_0}\left(\sin\varphi - \varphi\cos\varphi_0\right) = \text{const}.$$ It remains to substitute the initial values: $$\dot\varphi_i = -\frac{nK\omega_0}{E_0}\,\mathcal E_i.$$

\par\smallskip\noindent\textbf{Answer:} $$\dot\varphi^2 = \left(\frac{nK\omega_0}{E_0}\,\mathcal E_i\right)^2 + \frac{nK\omega_0^2\,eU}{\pi E_0}\left(\sin\varphi_i - \varphi_i\cos\varphi_0\right) - \frac{nK\omega_0^2\,eU}{\pi E_0}\left(\sin\varphi - \varphi\cos\varphi_0\right).$$
\end{solcard}

\begin{solcard}{C5}{Determine for which values of $\varphi_i$ and $\mathcal E_i$ the particle is captured into the acceleration process, i.e.\ on average it gains the same amount of energy as the equilibrium one. Qualitatively sketch the boundaries of the obtained region in the $(\varphi_i,\mathcal E_i)$ plane for $\varphi_0 = 0$ (resonant acceleration) and $\varphi_0 = \pi/4$.}{1.80}

Let us use the result of the previous point: $$\dot\varphi^2 + \frac{nK\omega_0^2\,eU}{\pi E_0}\left(\sin\varphi - \varphi\cos\varphi_0\right) = \text{const} = \left(\frac{nK\omega_0}{E_0}\,\mathcal E_i\right)^2 + \frac{nK\omega_0^2\,eU}{\pi E_0}\left(\sin\varphi_i - \varphi_i\cos\varphi_0\right).$$ Note that capture into acceleration is equivalent to bounded motion of $\varphi$: $$\dot\varphi^2 + V(\varphi) = \text{const} = H.$$ We obtain the problem of studying motion in some ``potential'' $V$ with ``energy'' $H$. Consider the case $K > 0$. Let us construct the qualitative shape of $V$ for this case.

\begin{figure}[H]
  \centering
  \includegraphics[width=0.72\textwidth]{T3_S_EN_assets/fig1.jpg}
\end{figure}

The red line indicates the maximum level of ``energy'' at which the motion is finite (occurs in a bounded region). For higher ``energy'' $\varphi$ can grow without bound (motion to the right on the graph). Let us find this value of $H$. It is determined by the value of the maximum of $V$: $$\frac{dV}{d\varphi} = \frac{nK\omega_0^2\,eU}{\pi E_0}\left(\cos\varphi - \cos\varphi_0\right).$$ The maximum corresponds to $\varphi = |\varphi_0|$: $$H = \frac{nK\omega_0^2\,eU}{\pi E_0}\left(\sin|\varphi_0| - |\varphi_0|\cos\varphi_0\right).$$ Then the condition on $\varphi_i$ and $\mathcal E_i$ takes the form $$\left(\frac{nK\omega_0}{E_0}\,\mathcal E_i\right)^2 + \frac{nK\omega_0^2\,eU}{\pi E_0}\left(\sin\varphi_i - \varphi_i\cos\varphi_0\right) \leqslant \frac{nK\omega_0^2\,eU}{\pi E_0}\left(\sin|\varphi_0| - |\varphi_0|\cos\varphi_0\right).$$ It is worth noting that there is also an additional condition on $\varphi_i$: initially the particle must be located in the ``potential well,'' which corresponds to $\varphi_i \leqslant |\varphi_0|$.

Note that for $K < 0$ the situation is analogous, but now the potential has a maximum at $\varphi = -|\varphi_0|$. Then the maximum ``energy'' is $$H = \frac{nK\omega_0^2\,eU}{\pi E_0}\left(-\sin|\varphi_0| + |\varphi_0|\cos\varphi_0\right),$$ with the analogous restriction on the phase $\varphi_i \geqslant -|\varphi_0|$. Note that we can generalise to arbitrary sign of $K$: $$H = \frac{n|K|\omega_0^2\,eU}{\pi E_0}\left(\sin|\varphi_0| - |\varphi_0|\cos\varphi_0\right),$$ $$\frac{K}{|K|}\,\varphi_i \leqslant |\varphi_0|.$$ The final expression for the desired region is:

\par\smallskip\noindent\textbf{Answer:} $$\left(\frac{nK\omega_0}{E_0}\,\mathcal E_i\right)^2 + \frac{nK\omega_0^2\,eU}{\pi E_0}\left(\sin\varphi_i - \varphi_i\cos\varphi_0\right) \leqslant \frac{n|K|\omega_0^2\,eU}{\pi E_0}\left(\sin|\varphi_0| - |\varphi_0|\cos\varphi_0\right);$$ $$\frac{K}{|K|}\,\varphi_i \leqslant |\varphi_0|.$$ The boundary of the region is given by $$\left(\frac{nK\omega_0}{E_0}\,\mathcal E_i\right)^2 + \frac{nK\omega_0^2\,eU}{\pi E_0}\left(\sin\varphi_i - \varphi_i\cos\varphi_0\right) = \frac{n|K|\omega_0^2\,eU}{\pi E_0}\left(\sin|\varphi_0| - |\varphi_0|\cos\varphi_0\right);$$ $$\frac{K}{|K|}\,\varphi_i \leqslant |\varphi_0|.$$

\medskip
\noindent\textbf{Case $\varphi_0 = 0$:} $$\left(\frac{nK\omega_0}{E_0}\,\mathcal E_i\right)^2 + \frac{nK\omega_0^2\,eU}{\pi E_0}\left(\sin\varphi_i - \varphi_i\right) = 0;\qquad K\,\varphi_i \leqslant 0.$$ It is easy to note that for $\varphi_i \neq 0$: $$K\left(\sin\varphi_i - \varphi_i\right) = -K\varphi_i\!\left(1-\frac{\sin\varphi_i}{\varphi_i}\right).$$ Since $1-\frac{\sin\varphi_i}{\varphi_i} > 0$, we have $$K\left(\sin\varphi_i - \varphi_i\right) = -K\varphi_i\!\left(1-\frac{\sin\varphi_i}{\varphi_i}\right) > 0.$$ Substituting into the first inequality, we obtain $$\left(\frac{nK\omega_0}{E_0}\,\mathcal E_i\right)^2 < 0,$$ so the only point satisfying the conditions is $(\varphi_i,\mathcal E_i) = (0,0)$. It is worth noting that this fact is obvious from consideration of the energy analogy: for $\varphi_0 = 0$ the potential has no minimum, only an inflection point at zero, so the only case of bounded motion is equilibrium at zero.

\begin{figure}[H]
  \centering
  \includegraphics[width=0.72\textwidth]{T3_S_EN_assets/fig2.jpg}
\end{figure}

\medskip
\noindent\textbf{Case $\varphi_0 = \pi/4$:} $$\left(\frac{nK\omega_0}{E_0}\,\mathcal E_i\right)^2 + \frac{nK\omega_0^2\,eU}{\pi E_0}\!\left(\sin\varphi_i - \frac{\varphi_i}{\sqrt 2}\right) = \frac{n|K|\omega_0^2\,eU}{\pi E_0}\!\left(\frac{1}{\sqrt 2} - \frac{\pi}{4\sqrt 2}\right);$$ $$\frac{K}{|K|}\,\varphi_i \leqslant \frac{\pi}{4}.$$ We obtain two cases depending on the sign of $K$:

\begin{figure}[H]
  \centering
  \includegraphics[width=0.72\textwidth]{T3_S_EN_assets/fig3.jpg}
\end{figure}
\end{solcard}

\end{document}
